\section*{Chapter \ref{ch:g10:molecules-of-life} --- Molecules of Life}

\begin{solution}{exo:g10:molecules-of-life:1}
Oxygen is more electronegative than hydrogen, so $\mathrm{H_2O}$ has
partial charges ($\delta^-$ on O, $\delta^+$ on H). Water's poles attract
$\mathrm{Na^+}$ and $\mathrm{Cl^-}$, surrounding each ion with a hydration
shell and separating the crystal lattice.
\end{solution}

\begin{solution}{exo:g10:molecules-of-life:2}
Dehydration (condensation) joins \omterm{def:g10:molecules-of-life:polymer}{monomers} and releases water; \omterm{def:g10:molecules-of-life:polymer}{hydrolysis}
splits \omterm{def:g10:molecules-of-life:polymer}{polymers} by adding water. Gut digestion of starch is \omterm{def:g10:molecules-of-life:polymer}{hydrolysis}.
\end{solution}

\begin{solution}{exo:g10:molecules-of-life:3}
Examples: glucose (mono); sucrose or lactose (di); starch (plant energy
store) and cellulose (plant wall structure); glycogen also acceptable as
animal store.
\end{solution}

\begin{solution}{exo:g10:molecules-of-life:4}
Head (phosphate) \omterm{def:g10:molecules-of-life:phil-phob}{hydrophilic}; two fatty-acid tails \omterm{def:g10:molecules-of-life:phil-phob}{hydrophobic}. In water,
tails sequester away from water while heads face aqueous phases, so a
bilayer is the low-energy arrangement.
\end{solution}

\begin{solution}{exo:g10:molecules-of-life:5}
A \omterm{def:g10:molecules-of-life:proteins}{peptide bond} is the covalent link between the carboxyl group of one
\omterm{def:g10:molecules-of-life:proteins}{amino acid} and the amino group of the next (with loss of water). Five
\omterm{def:g10:molecules-of-life:proteins}{amino acids} in one chain form four \omterm{def:g10:molecules-of-life:proteins}{peptide bonds}.
\end{solution}

\begin{solution}{exo:g10:molecules-of-life:6}
Heat disrupts weak interactions holding secondary/tertiary structure;
active site shape is lost (denaturation). Cooling restores activity only
if refolding is possible; often aggregation or irreversible damage
prevents recovery.
\end{solution}

\begin{solution}{exo:g10:molecules-of-life:7}
All three: glucose \omterm{def:g10:molecules-of-life:polymer}{monomers}. Starch: plants, some branching, energy store.
Glycogen: animals, highly branched, short-term store. Cellulose: plants,
unbranched $\beta$ linkages, structural fibre.
\end{solution}

\begin{solution}{exo:g10:molecules-of-life:8}
A--T and C--G. If $\%$A $= 30$, then $\%$T $= 30$; remaining $40\%$ split
equally between G and C, so $\%$G $= \%$C $= 20$.
\end{solution}

\begin{solution}{exo:g10:molecules-of-life:9}
Amylase: \omterm{def:g10:molecules-of-life:proteins}{protein}. Olive oil: \omterm{def:g10:molecules-of-life:lipids}{lipid}. Cellulose: \omterm{def:g10:molecules-of-life:carbs}{carbohydrate}. Insulin:
\omterm{def:g10:molecules-of-life:proteins}{protein}. mRNA: \omterm{def:g10:molecules-of-life:nucleic}{nucleic acid}. Cholesterol: \omterm{def:g10:molecules-of-life:lipids}{lipid} (\omterm{def:g10:molecules-of-life:lipids}{steroid}). Glycogen:
\omterm{def:g10:molecules-of-life:carbs}{carbohydrate}.
\end{solution}

\begin{solution}{exo:g10:molecules-of-life:10}
\omterm{def:g10:molecules-of-life:lipids}{Fats} are highly reduced hydrocarbons with little oxygen; oxidation yields
more energy per gram than hydrated \omterm{def:g10:molecules-of-life:carbs}{carbohydrates}. Consequence: efficient,
lightweight long-term energy stores (and insulation) in animals.
\end{solution}

\begin{solution}{exo:g10:molecules-of-life:11}
Starch: iodine/KI --- blue-black. Reducing sugar: Benedict's (heat) ---
brick red/orange precipitate. \omterm{def:g10:molecules-of-life:proteins}{Protein}: biuret --- purple. Run on separate
samples; include positive and negative controls.
\end{solution}

\begin{solution}{pb:g10:molecules-of-life:1}
\textbf{Part I.} Bread: starch (carb), some \omterm{def:g10:molecules-of-life:proteins}{protein}, fibre (cellulose).
Oil: \omterm{def:g10:molecules-of-life:lipids}{lipids}. Fish: \omterm{def:g10:molecules-of-life:proteins}{protein}, some \omterm{def:g10:molecules-of-life:lipids}{lipid}, \omterm{def:g10:molecules-of-life:nucleic}{nucleic acids}. Fruit: sugars,
fibre, water. Milk: lactose, \omterm{def:g10:molecules-of-life:proteins}{protein} (casein), \omterm{def:g10:molecules-of-life:lipids}{lipid}. Storage
\omterm{def:g10:molecules-of-life:carbs}{polysaccharide}: starch (bread); structural: cellulose (plant \omterm{def:g10:the-cell:support}{cell walls}
in salad/fruit).

\textbf{Part II.} Starch $\to$ glucose (and maltose intermediates);
\omterm{def:g10:molecules-of-life:proteins}{protein} $\to$ \omterm{def:g10:molecules-of-life:proteins}{amino acids}; \omterm{def:g10:molecules-of-life:lipids}{fat} $\to$ glycerol + fatty acids. All digestion:
\omterm{def:g10:molecules-of-life:polymer}{hydrolysis}.

\textbf{Part III.} \omterm{def:g10:molecules-of-life:proteins}{Amino acids} $\to$ new \omterm{def:g10:molecules-of-life:proteins}{proteins}; sugars/fatty acids
$\to$ respiration or storage; fatty acids/glycerol $\to$ \omterm{def:g10:molecules-of-life:lipids}{triglycerides} or
membrane \omterm{def:g10:molecules-of-life:lipids}{lipids}; phosphates/bases recycled. Dietary \omterm{def:g10:molecules-of-life:nucleic}{DNA} is hydrolysed to
\omterm{def:g10:molecules-of-life:nucleic}{nucleotides}; it is not intact human genetic instruction --- \omterm{def:g10:molecules-of-life:proteins}{sequence} and
cellular context matter, and foreign \omterm{def:g10:molecules-of-life:nucleic}{DNA} is not installed as your genome
by eating.
\end{solution}
